% Article VII. Additive incorporation, 10-09-2026.
% RESULTADO_RECUPERADO: fuentes_conservadas/18_clausura_energetica.tex,
% lines 357--485: Gramian, weight, minimum extension, and Schur elimination;
% lines 149--180 and 235--295: normalization and uniqueness of the response.
% Previous residence in VII: 30_pantalla_y_respuesta.tex,
% vii:prop:incidencia, vii:eq:lambda-pantalla, vii:prop:minima-extension.
% FORMALIZACION_NUEVA / CERTIFICADO_NUEVO: differential of the minimum with
% variable incidence, boundary, and normalization; block-transported
% incidence and composition with connection coordinates.
% The family of routes realizing each incidence remains explicit.
% Neither that full differentiable family nor the evaluation of a matter
% spin amplitude is attributed to the historical owner.
% Dependencies: 30, 30c, and the variational convention of 31. This does
% not modify those fragments or by itself constitute a matter realization.

\section{Differential of the minimum boundary extension}
\label{vii:sec:corriente-extension-minima}

Cubic incidence simultaneously determines the minimum-extension field
and the weight of its response. For a boundary datum, both are computed
through the inverse incidence Gramian. We now develop the differential
of that same operation, retaining variations of the boundary,
transport, and normalization scale. In this way, the field entering the
response is selected by the screen variational problem.

\subsection{Minimizing field and response weight}
\label{vii:subsec:extension-minima-parametrizada}

Let \(H\) and \(Q\) be finite-dimensional Hermitian spaces, with
fixed inner products antilinear in the first variable. The real case
is obtained by replacing adjoints with transposes. On an open set of
real parameters \(\Omega\), consider \(C^1\) collections
\[
 \mathsf C(\theta):H\longrightarrow Q,\qquad
 b(\theta)\in Q,\qquad \chi(\theta)>0,
 \qquad \operatorname{Ran}\mathsf C(\theta)=Q.
\]
Surjectivity is equivalent to full row rank in orthonormal bases.
Define the Gramian, the conjugate potential, and the extension
\begin{equation}
 \mathsf L=\mathsf C\mathsf C^*,\qquad
 y=\mathsf L^{-1}b,\qquad f_{\min}=\mathsf C^*y,
 \qquad \Lambda=\chi\mathsf L^{-1}.
 \label{vii:eq:gramiano-extension-minima}
\end{equation}
The incidence \(\mathsf C\) and the cellular pairing \(M\) of
\eqref{vii:eq:coordenadas-corriente-material} are different operators.

\begin{proposition}[Variational selection of the extension]
\label{vii:prop:minimo-parametrizado}
The extension \(f_{\min}\) uniquely minimizes the norm among
fields \(f\in H\) satisfying \(\mathsf C f=b\). Its energy is
\begin{equation}
 E(\theta)=\chi\min_{\mathsf C f=b}\|f\|^2
 =\chi\langle b,\mathsf L^{-1}b\rangle
 =\langle b,\Lambda b\rangle.
 \label{vii:eq:energia-extension-minima}
\end{equation}
The field \(f_{\min}\), potential \(y\), and energy \(E\)
depend \(C^1\)-smoothly on the parameters.
\end{proposition}

\begin{proof}
Surjectivity of \(\mathsf C\) makes \(\mathsf C^*\) injective,
so \(\langle z,\mathsf Lz\rangle=\|\mathsf C^*z\|^2>0\)
for \(z\ne0\). Thus \(\mathsf L\) is invertible and
\(\mathsf C f_{\min}=\mathsf L y=b\). Every solution is
\(f_{\min}+z\), with \(z\in\ker\mathsf C\). Since
\(f_{\min}\in\operatorname{Ran}\mathsf C^*=(\ker\mathsf C)^\perp\),
\[
 \|f_{\min}+z\|^2=\|f_{\min}\|^2+\|z\|^2,
 \qquad
 \|f_{\min}\|^2=\langle y,\mathsf L y\rangle
 =\langle b,\mathsf L^{-1}b\rangle.
\]
These identities prove the minimum, its uniqueness, and its value.
Inversion is a differentiable operation on the open set of invertible
matrices; the expressions in \eqref{vii:eq:gramiano-extension-minima}
give the asserted regularity.
\end{proof}

The cubic realization of \ref{vii:prop:incidencia} corresponds to
\(H=\mathbb R^{\{-1,1\}^3}\), the screen quotient
\(Q=P_\square\mathbb R^6\), and the incidence \(\mathsf C=M:H\to Q\).
Its Gramian and normalization are
\begin{equation}
 \mathsf L_\square=12P_u+4P_-,\qquad
 \chi_m=\frac{112}{81\ell_m^2},\qquad
 \Lambda_m=\frac{28}{243\ell_m^2}(P_u+3P_-).
 \label{vii:eq:normalizacion-minimo-cubico}
\end{equation}
The factor \(112/81\) comes from soldering and the quadratic
compression defect of \eqref{vii:eq:lambda-pantalla}. If \(b\) has
the dimension of length, so do \(y\) and \(f_{\min}\); incidence
is dimensionless and \(E\) is dimensionless. Conversion to action
uses the normalization in the action line specified below.

\Needspace{15\baselineskip}
\subsection{Full differential of the minimum}
\label{vii:subsec:diferencial-minimo}

\begin{theorem}[Response to variations of incidence and boundary]
\label{vii:thm:diferencial-minimo}
For every real tangent direction \(\delta\theta\),
\begin{equation}
 \delta E
 =\delta\chi\,\langle b,y\rangle
   +2\chi\operatorname{Re}
       \langle y,\delta b-(\delta\mathsf C)f_{\min}\rangle.
 \label{vii:eq:diferencial-minimo-completo}
\end{equation}
If \(\chi=112/(81\ell_m^2)\), with variable \(\ell_m>0\),
\begin{equation}
 \delta E=-2E\frac{\delta\ell_m}{\ell_m}
 +2\chi_m\operatorname{Re}
       \langle y,\delta b-(\delta\mathsf C)f_{\min}\rangle.
 \label{vii:eq:diferencial-minimo-escala}
\end{equation}
\end{theorem}

\begin{proof}
Differentiating \(\mathsf L\mathsf L^{-1}=I_Q\) gives
\[
 \delta\mathsf L^{-1}
 =-\mathsf L^{-1}(\delta\mathsf L)\mathsf L^{-1},\qquad
 \delta\mathsf L=(\delta\mathsf C)\mathsf C^*
                      +\mathsf C(\delta\mathsf C)^*.
\]
Applying the product rule to \eqref{vii:eq:energia-extension-minima}
and using self-adjointness of \(\mathsf L^{-1}\) gives
\[
 \delta E=\delta\chi\,\langle b,y\rangle
 +2\chi\operatorname{Re}\langle y,\delta b\rangle
 -\chi\langle y,(\delta\mathsf L)y\rangle.
\]
The two terms in the last inner product are conjugates of one another,
and \(\mathsf C^*y=f_{\min}\); hence
\[
 \langle y,(\delta\mathsf L)y\rangle
 =2\operatorname{Re}\langle y,(\delta\mathsf C)f_{\min}\rangle.
\]
This yields \eqref{vii:eq:diferencial-minimo-completo}.
Finally, \(\delta\chi_m=-2\chi_m\delta\ell_m/\ell_m\)
gives \eqref{vii:eq:diferencial-minimo-escala}.
\end{proof}

The formula incorporates variation of the minimizing field through
the equation determining it. This can also be checked directly:
\(\mathsf C\delta f_{\min}=\delta b-(\delta\mathsf C)f_{\min}\)
and \(f_{\min}=\mathsf C^*y\) imply
\[
 \delta\|f_{\min}\|^2
 =2\operatorname{Re}\langle f_{\min},\delta f_{\min}\rangle
 =2\operatorname{Re}
     \langle y,\delta b-(\delta\mathsf C)f_{\min}\rangle.
\]
For example, if \(\mathsf C\) is fixed and \(b=\ell_m b_0\),
with fixed \(b_0\), the two terms of
\eqref{vii:eq:diferencial-minimo-escala} cancel and the energy
remains constant. This check jointly preserves the dimension of the
boundary datum and its normalization.

\subsection{Fixed quotient and changes of frame}
\label{vii:subsec:marcos-minimo}

The preceding derivatives are taken in fixed spaces \(H,Q\) with
fixed Hermitian inner products. In the six-face presentation, the
inverse always acts on \(Q=P_\square H_F\). For a collection of
quotients \(Q(\theta)\), first fix an isometric trivialization
\(V_Q(\theta):Q\to Q(\theta)\), and similarly for the domain.
The derivatives of these maps are included when differentiating the
expressions transported to the fixed spaces. The formulas apply on
each open set where rank is preserved. Changes in the energy inner
product require differentiating that inner product as well; a
Lorentzian metric and the positive inner product used in minimization
retain their distinct types here.

\begin{proposition}[Covariance of the minimum]
\label{vii:prop:covariancia-minimo}
Let \(g_H,g_Q\) be unitary. Under
\[
 \mathsf C'=g_Q\mathsf C g_H^*,\qquad b'=g_Qb,\qquad \chi'=\chi,
\]
one has \(y'=g_Qy\), \(f'_{\min}=g_Hf_{\min}\), and \(E'=E\).
When changes of frame depend on the parameter, the full differential
preserves this equality. In particular, a pure frame variation has
\(\delta E=0\).
\end{proposition}

\begin{proof}
The transformed Gramian is \(\mathsf L'=g_Q\mathsf Lg_Q^*\),
which gives the transformations of \(y,f_{\min}\) and equality
of energies. For a pure frame variation, at a point with identity
frames, let \(A_Q^*=-A_Q\) and \(A_H^*=-A_H\) be their
generators. Then
\[
 \delta\mathsf C=A_Q\mathsf C-\mathsf C A_H,\qquad
 \delta b=A_Qb,\qquad
 \delta b-(\delta\mathsf C)f_{\min}=\mathsf C A_Hf_{\min}.
\]
Its contribution to \eqref{vii:eq:diferencial-minimo-completo} is
\(2\chi\operatorname{Re}\langle f_{\min},A_Hf_{\min}\rangle=0\).
The proof thus includes the terms arising from moving frames.
\end{proof}

\subsection{Transported incidence and connection components}
\label{vii:subsec:incidencia-conexion-minima}

Connection dependence is specified before contraction is performed.
An explicit block realization of incidence is obtained as follows.
Let \(\mathcal K\) be a common Hermitian fibre,
\(F=\{(i,\epsilon):1\leq i\leq3,\ \epsilon=\pm1\}\), and
\(V=\{-1,1\}^3\). For each incidence \(\nu_i=\epsilon\),
the realization being used declares a route \(\gamma_{i,\epsilon,\nu}\)
from the vertex fibre to the face fibre, with transport
\(U_{i,\epsilon,\nu}\). The trivializations identify these
fibres with \(\mathcal K\). Define
\begin{equation}
 \begin{split}
 (\mathsf B_Uf)_{i,\epsilon}
   &=\sum_{\nu_i=\epsilon}U_{i,\epsilon,\nu}f_\nu,\\
 \mathsf C_U&=(q_\square\otimes I_{\mathcal K})\mathsf B_U:
   \bigoplus_{\nu\in V}\mathcal K\longrightarrow Q\otimes\mathcal K,
 \qquad q_\square=P_\square:H_F\to Q.
 \end{split}
 \label{vii:eq:incidencia-transportada-minima}
\end{equation}
This is an explicit variational continuation of cubic incidence in
the declared route representation. For identity transports,
\(\mathsf C_U=M\otimes I_{\mathcal K}\), and
\eqref{vii:eq:normalizacion-minimo-cubico}, tensored with the fibre,
is recovered. The Gramian remains positive in a neighbourhood of that
realization, since positive definiteness is an open condition. On each
full-rank domain, the weight and field are fixed by
\eqref{vii:eq:gramiano-extension-minima}. The routes realizing the
incidences, their orientations, and their memories are data of this
realization of APP--TRIT--TPK transport and accompany the construction.

If \(\omega=(\omega_j)\) are real connection coordinates,
the derivative of \(\mathsf C_U\) is computed by substituting
in each block
\begin{equation}
 \partial_j U_\gamma
 =\sum_{k=1}^{N}
    U_N\cdots U_{k+1}(\partial_jU_k)U_{k-1}\cdots U_1,
 \qquad U_\gamma=U_N\cdots U_1.
 \label{vii:eq:derivada-incidencia-rutas}
\end{equation}
Empty products are identities. The formula follows from the product
rule and preserves the position of each operation along the route.
Thus \(\mathsf C_j:=\partial_j\mathsf C_U\) is obtained from
a finite sum of matrices known in the chosen realization, followed
by the projection \(q_\square\otimes I\). If a moving quotient
is used, its derivative is incorporated as in
\ref{vii:subsec:marcos-minimo}. The fixed integer cubic matrix
corresponds to the case \(\mathsf C_j=0\); the connection response
comes from the transport dependence actually declared in
\eqref{vii:eq:incidencia-transportada-minima}.

\begin{theorem}[Variational components of the minimum extension]
\label{vii:thm:componentes-minimo-conexion}
For the screen action
\(S_{\min}=\mathfrak a_m E\), where \(\mathfrak a_m\) is its
normalization in the action line, write
\[
 b_j=\partial_jb,\qquad \chi_j=\partial_j\chi,
 \qquad a_j=\partial_j\mathfrak a_m.
\]
With the convention \(\delta S_{\min}=-\tfrac12\sum_j s_j\delta\omega_j\),
the components are
\begin{equation}
 \begin{split}
 s_j={}&4\mathfrak a_m\chi\operatorname{Re}
             \langle y,\mathsf C_jf_{\min}-b_j\rangle\\
      &-2\mathfrak a_m\chi_j\langle b,y\rangle-2a_jE.
 \end{split}
 \label{vii:eq:componentes-minimo-conexion}
\end{equation}
For the cubic normalization, the second term is equivalent to
\(4\mathfrak a_m E\,\partial_j\ell_m/\ell_m\).
If the cellular pairing of this same action has invertible matrix
\(\mathsf M_{\rm cel}\), its current coordinates are
\(\sigma=\mathsf M_{\rm cel}^{-1}s\).
\end{theorem}

\begin{proof}
The product rule gives
\(\partial_jS_{\min}=a_jE+\mathfrak a_m\partial_jE\).
Substituting \eqref{vii:eq:diferencial-minimo-completo} and
multiplying by \(-2\) proves \eqref{vii:eq:componentes-minimo-conexion}.
The derivative of \(\chi_m\) gives the scale expression.
Finally, comparing
\(\delta S_{\min}=-\tfrac12\delta\omega^{\mathsf T}s\)
with \(-\tfrac12\delta\omega^{\mathsf T}\mathsf M_{\rm cel}\sigma\)
for every variation gives \(\mathsf M_{\rm cel}\sigma=s\),
and its inverse determines the coordinates.
\end{proof}

The construction recovers a field and a weight jointly selected by
incidence and evaluates their response to variations of the same
realization. The boundary datum \(b\), the route configuration, the scale
\(\ell_m\), the normalization \(\mathfrak a_m\), and the
cellular pairing retain their representation rules. Interpretation as a
matter current uses these objects and the full action. The Cartan
response and its subsequent metric density are composed with it in
the corresponding matter domain, retaining any dependence of the
current on the connection.
